% ===========================================================================
%  6 — Arquitetura de nuvem e implantacao
% ===========================================================================
\chapter{Arquitetura de nuvem e implantação}
\label{cap:arqnuvem}

Este capítulo descreve \emph{onde} o \emph{software} do
Capítulo~\ref{cap:arqsw} roda: quais recursos gerenciados existem, como o
tráfego atravessa a plataforma, como a identidade de carga de trabalho
substitui chaves e onde estão as fronteiras de rede.

A infraestrutura é provisionada por Terraform~\cite{brikman2022terraform} num
\textbf{repositório separado}. Este documento descreve o \emph{contrato} entre os
dois repositórios --- nomes de recursos esperados, variáveis consumidas,
artefatos publicados --- e não o código Terraform.

\begin{regra}
\textbf{Este repositório nunca cria recurso GCP imperativamente.} Nenhum
\emph{script} de \arq{scripts/} chama \cd{gcloud ... create}. O CI apenas
\emph{usa} recursos que já existem: envia imagem para um Artifact Registry
existente, executa um Cloud Run Job existente, sincroniza um \emph{bucket}
existente. Recurso criado por \emph{script} é recurso que o Terraform não sabe
que existe --- e que ninguém destrói quando o contrato termina.
\end{regra}

\section{Princípios de nuvem}

\begin{description}
  \item[Uma \emph{stack} por contrato.] Todos os nomes de recurso derivam de
    \cd{\{contrato\} = \{app\_name\}-\{env\}}. No piloto,
    \cd{plataforma-unooria-dev}. Replicar um contrato é instanciar o
    \emph{module} Terraform com outro prefixo.
  \item[Serviço gerenciado antes de orquestração.] Cloud Run em vez de GKE;
    Cloud SQL em vez de PostgreSQL autogerido; Pub/Sub em vez de \emph{broker}
    próprio. O critério não é moda: é quantas peças uma pessoa consegue operar de
    plantão às oito da manhã.
  \item[Sem chave de longa duração.] O CI autentica por Workload Identity
    Federation; a API autentica no Vertex por ADC. Nenhum arquivo JSON de conta
    de serviço existe em nenhum lugar do fluxo.
  \item[Reversão por ponteiro.] Três níveis: revisão do Cloud Run (hub e API),
    \cd{versaoAtiva} no manifesto (experiências) e versionamento de objetos
    (planos de aula). Nenhum deles exige \emph{build}.
  \item[Estático na borda, dinâmico no mínimo.] O que pode ser servido pela CDN
    é servido pela CDN. O que sobra para o Cloud Run é regra de negócio.
\end{description}

\section{C4 nível 1 de implantação --- o contrato}

\begin{figure}[H]
\centering
\fitwidth{%
\begin{tikzpicture}[x=1cm, y=1cm,
  caixa/.style={text width=28mm, minimum height=14mm},
  ia/.style={servico, fill=unooLavender, draw=unooLavenderSoft},
  zonaia/.style={regiao, draw=unooLavender, fill=unooLavenderWash},
  zonavpc/.style={regiao, draw=unooGreen!60, fill=unooGreenWash},
  zonalbl/.style={c4boundarylbl, fill=unooSurface, anchor=south west}]

  % grade: colunas A = 1.3 (fora do Google), B = 5.7, C = 10.4, D = 14.6;
  % linhas 0, -2.8, -5.8, -8.6. As rotas correm nos corredores entre elas.

  % ------------------------------------------------ fora da stack
  \node[internet, minimum width=27mm] (net) at (1.3,0) {Navegador\\\emph{da escola}};
  \node[externo, caixa] (idp) at (1.3,-2.8)
    {\textbf{WSO2 Identity}\\\textbf{Server}\desc{projeto próprio, vários contratos}};
  \node[externo, caixa] (lav) at (1.3,-8.6)
    {\textbf{LiveAvatar (HeyGen)}\desc{vídeo da Ajuda; fora do Google}};

  % ------------------------------------------------ borda
  \node[servicoger, text width=73mm, minimum height=13mm] (lb) at (8.05,0)
    {\textbf{Cloud Load Balancing}\desc{global, um IP, \emph{host rules}: hub, API, \emph{bundles}, \emph{assets}}};
  \node[servico, caixa, fill=unooGreen!75, draw=unooGreen!55!black] (gcs) at (14.6,0)
    {\textbf{Cloud Storage}\desc{\texttt{-exp}, \texttt{-assets}: CDN\\\texttt{-planos}: privado}};

  % ------------------------------------------------ aplicação
  \node[servico, caixa, text width=30mm] (hub) at (5.7,-2.8)
    {\textbf{Cloud Run}\\\texttt{\{contrato\}-frontend}\desc{hub (Next.js)}};
  \node[servico, caixa, minimum height=16mm] (api) at (10.4,-2.8)
    {\textbf{Cloud Run}\\\texttt{\{contrato\}-api}\desc{API (NestJS)}};
  \node[servicobroker, caixa] (ps) at (14.6,-2.8)
    {\textbf{Pub/Sub}\desc{telemetria analítica}};

  % ------------------------------------------------ dados (VPC)
  \node[volume, text width=25mm, minimum height=15mm] (sql) at (5.7,-5.8)
    {\textbf{Cloud SQL}\\\texttt{\{contrato\}-pg}\desc{HA, PITR, \emph{pooler}}};

  % ------------------------------------------------ IA gerenciada
  \node[ia, caixa] (agente) at (10.4,-5.8)
    {\textbf{Agent Engine}\\\texttt{agente-professor}\desc{Python + ADK}};
  \node[ia, caixa] (gemini) at (14.6,-5.8)
    {\textbf{Gemini}\desc{\texttt{global}; o Live em \texttt{us-central1}}};
  \node[ia, caixa] (busca) at (10.4,-8.6)
    {\textbf{Vertex AI Search}\desc{dois \emph{data stores}\\em \texttt{global}}};
  \node[servico, caixa, fill=unooGreen!75, draw=unooGreen!55!black] (bag) at (14.6,-8.6)
    {\textbf{Cloud Storage}\\\texttt{-agente-professor}\desc{BNCC e acervo; privado}};

  % ------------------------------------------------ zonas (atrás dos nós)
  \begin{scope}[on background layer]
    % a nuvem primeiro: o que vem depois no mesmo plano fica por cima dela
    \node[nuvem, fit={(lb) (gcs) (bag) (hub.west |- bag.south)},
          inner xsep=16pt, inner ysep=26pt, yshift=-4pt] (gcp) {};
    \node[zonavpc, fit=(sql), inner sep=6pt] (vpc) {};
    \node[zonaia, fit=(agente) (gemini) (busca) (bag), inner sep=7pt] (zia) {};
  \end{scope}
  \node[nuvemlbl, anchor=south west] at ([xshift=6pt, yshift=-1pt]gcp.north west)
    {Google Cloud --- projeto do contrato};
  \node[zonalbl, anchor=north west, text=unooGreen!60!black]
    at ([yshift=-1pt]vpc.south west) {VPC (só IP privado)};
  \node[zonalbl, anchor=north west, text=unooLavenderSoft]
    at ([yshift=-1pt]zia.south west) {IA gerenciada --- Vertex AI};

  % ------------------------------------------------ arestas
  \draw[netpub] (net.east) -- (lb.west) node[portlbl, midway, above=1pt]{443};
  \draw[netpub] (lb.east) -- (gcs.west);
  \draw[netpub] (lb.south -| hub.north) -- (hub.north);
  \draw[netpub] (lb.south -| api.north) -- (api.north);
  \draw[netpub] ([yshift=3mm]hub.east) -- ([yshift=3mm]hub.east -| api.west)
    node[portlbl, midway, above=1pt]{BFF};
  \draw[netpub] (hub.west) -- (idp.east) node[portlbl, pos=0.3, above=1pt]{OIDC};
  \draw[netpub] ([xshift=-3mm]api.north east) -- ++(0,0.65) -| (gcs.south)
    node[portlbl, pos=0.25, above=1pt]{\emph{assets}, planos};
  \draw[netbi] (api.east) -- (ps.west);
  \draw[netpub] ([xshift=-3mm]api.south east) |- (12.5,-4.45) -| (gemini.north)
    node[portlbl, pos=0.25, above=1pt]{HTTPS, WSS};
  \draw[netpub] (api.south) -- (agente.north) node[portlbl, pos=0.4, right=2pt]{SSE};
  \draw[netpub] (agente.east) -- (gemini.west);
  \draw[netpub] (agente.south) -- (busca.north) node[portlbl, midway, right=2pt]{busca};
  \draw[netpriv] (bag.west) -- (busca.east) node[portlbl, midway, above=1pt]{lê};
  % LiveAvatar: a API abre a sala; o vendor chama /ajuda/llm de volta
  \draw[netbi] ([yshift=-5mm]api.west) -- (8.35,-3.3) -- (8.35,-8.6) -- (lav.east)
    node[portlbl, pos=0.55, above=1pt]{sessões; \texttt{/ajuda/llm}};
  \draw[netpriv] ([yshift=-1.5mm]api.west) -- (7.75,-2.95) |- (sql.east)
    node[portlbl, pos=0.25, left=1pt]{5432};
  % vídeo direto ao navegador
  \draw[netpub] (net.west) -- (-0.75,0) |- (lav.west)
    node[portlbl, pos=0.25, sloped, above=1pt]{WebRTC};

\end{tikzpicture}}
\caption[Implantação de um contrato]{Implantação de um contrato. Um único
  \emph{Load Balancer} global recebe todo o tráfego HTTPS e distribui por
  \emph{host rules}: hub, API, \emph{bundles} e \emph{assets} respondem em
  domínios diferentes com o \emph{mesmo} IP. O Cloud SQL só tem IP privado ---
  daí a existência do Cloud Run Job de migração dentro da VPC
  (Figura~\ref{fig:deploy-detalhe}). Só a API fala com a IA: com o Gemini
  direto (a geração, e o Gemini Live do mascote da Ajuda, \adrr{ADR-158}, por
  WebSocket) e com o agente de IA do professor no Agent Engine, por SSE
  (\adrr{ADR-126}); o agente, por sua vez, usa o mesmo Gemini e os dois
  \emph{data stores} do Vertex AI Search, que leem o \emph{bucket} privado do
  agente. O navegador nunca alcança a zona de IA. O LiveAvatar, da Ajuda com
  avatar (\adrr{ADR-132}), é o único externo com tráfego nos dois sentidos: a
  API abre e fecha a sala em \texttt{api.liveavatar.com}, e o \emph{vendor}
  entra pelo mesmo LB para chamar \texttt{/ajuda/llm}; o vídeo vai dele direto
  ao navegador da escola.}
\label{fig:deploy-contrato}
\end{figure}

\begin{legendadiag}
\textcolor{unooOrange}{\textbf{---}} tráfego público (TLS) \quad
\textcolor{unooGreen!70!black}{\textbf{---}} tráfego privado \quad
\textcolor{unooInkSoft}{\textbf{---}} nos dois sentidos \quad
\amostra{servicoger} borda \quad
\amostra{servico} aplicação \quad
\amostra{volume} dados \quad
\amostra{servicobroker} mensageria \quad
\amostra{servico, fill=unooLavender, draw=unooLavenderSoft} IA gerenciada \quad
\amostra{externo} fora da \emph{stack}
\end{legendadiag}

\section{Diagrama de implantação detalhado}

A Figura~\ref{fig:deploy-detalhe} abre o mesmo desenho no nível em que ele é
operado, em execução: a borda com o WAF, os \emph{buckets} separados por regime
de acesso, o conector da VPC, o \emph{job} de migração, a mensageria e a região
de cada peça de IA. As identidades de máquina estão na
Figura~\ref{fig:identidade-maquina}.

\begin{figure}[H]
\centering
\fitpage{%
\begin{tikzpicture}[x=1cm, y=1cm,
  caixa/.style={text width=27mm, minimum height=14mm},
  ia/.style={servico, fill=unooLavender, draw=unooLavenderSoft},
  balde/.style={servico, fill=unooGreen!75, draw=unooGreen!55!black},
  zonaia/.style={regiao, draw=unooLavender, fill=unooLavenderWash},
  zonavpc/.style={regiao, draw=unooGreen!60, fill=unooGreenWash},
  zonalbl/.style={c4boundarylbl, fill=unooSurface, anchor=north west},
  importa/.style={-{Stealth[length=5pt,width=4pt]}, line width=0.9pt,
                  draw=unooInkSoft, dashed, dash pattern=on 3pt off 2pt}]

  % grade: externos em x = 0.6; colunas c1 = 4.4, c2 = 7.8, c3 = 11.2,
  % c4 = 15.0; linhas 0, -3, -6, -9. Rotas nos corredores y = -1.5,
  % -4.6 e -7.6.

  % ------------------------------------------------ fora do Google
  \node[internet, minimum width=25mm] (net) at (0.6,0) {Navegador};
  \node[externo, text width=24mm, minimum height=14mm] (idp) at (0.6,-3)
    {\textbf{WSO2 IS}\desc{projeto próprio}};

  % ------------------------------------------------ borda global
  \node[servicoger, caixa] (armor) at (4.4,0)
    {\textbf{Cloud Armor}\desc{WAF nos \emph{backends}; SQLi, XSS, RCE\ldots}};
  \node[servicoger, caixa] (lb) at (7.8,0)
    {\textbf{Load Balancer}\desc{1 IP, \emph{host rules}; um certificado por domínio}};
  \node[servicoger, caixa] (cdn) at (11.2,0)
    {\textbf{Cloud CDN}\desc{\emph{bundles} e \emph{assets}}};

  % ------------------------------------------------ buckets
  \node[balde, text width=39mm, align=left, minimum height=22mm] (gcs) at (15.6,-0.45)
    {\textbf{Cloud Storage}\\[2pt]
     \texttt{-exp}: \emph{bundles}; público, CDN\\
     \texttt{-assets}: mídia; público, CDN\\
     \texttt{-planos}: \textbf{privado}, versionado\\
     \texttt{-exp-mapas}: privado\\
     \texttt{-agente-professor}: privado};

  % ------------------------------------------------ aplicação (us-east1)
  \node[servico, caixa] (hub) at (4.4,-3)
    {\textbf{Cloud Run} \texttt{-frontend}\desc{hub; \texttt{min=1}}};
  \node[servico, caixa, text width=29mm] (api) at (7.8,-3)
    {\textbf{Cloud Run} \texttt{-api}\desc{CPU sempre alocada, \texttt{min}$\geq$1, conc.~80}};
  \node[ia, caixa] (agente) at (11.2,-3)
    {\textbf{Agent Engine}\\\texttt{agente-professor}\desc{\texttt{us-east1}; SA \texttt{-agente}}};

  % ------------------------------------------------ VPC
  \node[servicoger, caixa] (conn) at (4.4,-6)
    {\textbf{Serverless VPC Access}\desc{conector}};
  \node[volume, text width=25mm, minimum height=15mm] (sql) at (4.4,-9)
    {\textbf{Cloud SQL} \texttt{-pg}\desc{PostgreSQL HA, PITR, \emph{pooler}; sem IP público}};
  \node[servicoger, caixa] (job) at (7.8,-9)
    {\textbf{Cloud Run Job}\\\texttt{-migrate}\desc{\texttt{prisma migrate deploy}}};

  % ------------------------------------------------ IA global
  \node[ia, caixa] (gemini) at (11.2,-6)
    {\textbf{Gemini}\desc{\texttt{global}, reserva em \texttt{us}; Live em \texttt{us-central1}}};
  \node[ia, caixa] (busca) at (15.6,-6)
    {\textbf{Vertex AI Search}\desc{\texttt{-habilidades} e \texttt{-acervo}; \texttt{global}}};

  % ------------------------------------------------ mensageria
  \node[servicobroker, text width=31mm, minimum height=14mm] (ps) at (11.2,-9)
    {\textbf{Pub/Sub}\desc{\emph{topic} \texttt{-telemetria}\\\emph{push sub} \texttt{-telemetria-ingest}}};

  % ------------------------------------------------ zonas
  \begin{scope}[on background layer]
    \node[nuvem, fit={(armor) (gcs) (sql) (ps) (busca)},
          inner xsep=14pt, inner ysep=22pt, yshift=-3pt] (gcp) {};
    \node[regiao, fit=(armor) (lb) (cdn), inner sep=6pt] (borda) {};
    \node[zonavpc, fit=(conn) (sql), inner sep=6pt] (vpc) {};
    \node[zonaia, fit=(agente) (gemini) (busca), inner sep=7pt] (zia) {};
  \end{scope}
  \node[nuvemlbl, anchor=south west] at ([xshift=6pt, yshift=-1pt]gcp.north west)
    {Google Cloud --- projeto do contrato \quad \texttt{\{contrato\} = \{app\_name\}-\{env\}}};
  \node[zonalbl] at ([yshift=-1pt]borda.south west) {Borda global};
  \node[zonalbl, text=unooGreen!60!black] at ([yshift=-1pt]vpc.south west) {VPC};
  \node[zonalbl, text=unooLavenderSoft, anchor=north east]
    at ([yshift=-1pt]zia.south east) {IA gerenciada --- Vertex AI};

  % ------------------------------------------------ arestas
  \draw[netpub] (net.east) -- (armor.west) node[portlbl, midway, above=1pt]{443};
  \draw[netpub] (armor.east) -- (lb.west);
  \draw[netpub] (lb.east) -- (cdn.west);
  \draw[netpub] (cdn.east) -- (cdn.east -| gcs.west);
  \draw[netpub] ([xshift=-5mm]lb.south) |- (4.4,-1.5) -- (hub.north);
  \draw[netpub] ([xshift=5mm]lb.south) -- ([xshift=5mm]lb.south |- api.north);
  \draw[netpub] (hub.east) -- (api.west) node[portlbl, midway, above=1pt]{BFF};
  \draw[netpub] (hub.west) -- (idp.east) node[portlbl, pos=0.35, above=1pt]{OIDC};
  \draw[netpub] ([xshift=-4mm]api.north east) |- (gcs.west |- 0,-1.5)
    node[portlbl, pos=0.7, above=1pt]{\texttt{-assets}, \texttt{-planos}};
  \draw[netpub] (api.east) -- (agente.west) node[portlbl, midway, above=1pt]{SSE};
  \draw[netpub] (agente.south) -- (gemini.north);
  \draw[netpub] (agente.east) -| (busca.north) node[portlbl, pos=0.3, above=1pt]{busca};
  \draw[importa] ([xshift=8mm]busca.north |- gcs.south) -- ([xshift=8mm]busca.north)
    node[portlbl, pos=0.8, left=1pt]{lê \texttt{-agente-professor}};
  \draw[netpriv] ([xshift=-9mm]api.south) |- (4.4,-4.6) -- (conn.north);
  \draw[netpriv] (conn.south) -- (sql.north) node[portlbl, midway, right=2pt]{5432};
  \draw[netpriv] (job.north) |- (conn.east);
  \draw[netpub] ([xshift=9mm]api.south) |- (gemini.west)
    node[portlbl, pos=0.75, above=1pt]{HTTPS, WSS};
  \draw[netbi] ([xshift=4mm]api.south) |- (11.2,-7.85) -- (ps.north)
    node[portlbl, pos=0.35, above=1pt]{publica; \emph{push} \texttt{/telemetry/ingest}};

\end{tikzpicture}}
\caption[Implantação detalhada --- o caminho da requisição]{Implantação
  detalhada, em execução. Os \emph{buckets} têm regimes de acesso
  deliberadamente diferentes: o de \emph{bundles} e o de \emph{assets} são
  públicos atrás da CDN (caminho versionado, portanto imutável); o de planos de
  aula é \textbf{privado}, sem CDN, com versionamento de objetos, e só é lido
  pela API depois de verificar papel e posse; o de \emph{source maps} e o do
  agente são privados. O Cloud SQL não tem IP público: a API chega a ele pelo
  conector de VPC, e é isso que obriga a migração a rodar como \emph{Job} com o
  mesmo conector. O Agent Engine roda na região do contrato; o modelo e os
  \emph{data stores}, em \texttt{global}. O Cloud Armor aplica as regras do WAF a
  todo caminho; só a de injeção de cabeçalho (921150) é excluída nos dois do
  chat do agente (Capítulo~\ref{cap:tenancy}). A
  Ajuda com avatar, fora do Google, está na Figura~\ref{fig:deploy-contrato}; o
  que o CI publica, na Figura~\ref{fig:deploy-publicacao}.}
\label{fig:deploy-detalhe}
\end{figure}

\begin{legendadiag}
\textcolor{unooOrange}{\textbf{---}} público (TLS) \quad
\textcolor{unooGreen!70!black}{\textbf{---}} privado (VPC) \quad
\textcolor{unooInkSoft}{\textbf{---}} nos dois sentidos \quad
\textcolor{unooInkSoft}{\textbf{- - -}} importação interna do Google \quad
\amostra{servicoger} borda e infraestrutura \quad
\amostra{servico} aplicação \quad
\amostra{servico, fill=unooGreen!75, draw=unooGreen!55!black} armazenamento \quad
\amostra{servico, fill=unooLavender, draw=unooLavenderSoft} IA gerenciada \quad
\amostra{servicobroker} mensageria
\end{legendadiag}
A Figura~\ref{fig:deploy-publicacao} mostra o outro lado: o que o CI escreve em
cada recurso, e com qual credencial. Todo passo usa o mesmo \emph{token} curto da
\cd{DEPLOY\_SA}, obtido pelo WIF sem chave JSON; a única exceção é a ativação da
versão no manifesto, que fala com a API como qualquer cliente, com um
\emph{token} RS256 de cinco minutos assinado pela
\cd{\{contrato\}-token-signing-key}.

\begin{figure}[H]
\centering
\fitwidth{%
\begin{tikzpicture}[x=1cm, y=1cm,
  alvo/.style={text width=38mm, minimum height=13mm},
  ia/.style={servico, alvo, fill=unooLavender, draw=unooLavenderSoft},
  balde/.style={servico, alvo, fill=unooGreen!75, draw=unooGreen!55!black},
  puxa/.style={-{Stealth[length=5pt,width=4pt]}, line width=0.8pt,
               draw=unooInkSoft, dashed, dash pattern=on 3pt off 2pt}]

  % passos do CI à esquerda; recursos numa coluna só, em x = 11.6;
  % o barramento do token corre em x = 7.4 e o da imagem em x = 14.2

  \node[externo, text width=26mm, minimum height=14mm] (gh) at (0.8,0)
    {\textbf{GitHub Actions}\desc{\texttt{deploy.yml}; \emph{job} com \texttt{id-token: write}}};
  \node[servicoger, text width=27mm, minimum height=14mm] (wif) at (4.6,0)
    {\textbf{WIF} \emph{pool} \texttt{github}\desc{troca o OIDC do GitHub pelo \emph{token} curto da \texttt{DEPLOY\_SA}}};

  \node[servicoger, alvo] (ar) at (11.8,4.0)
    {\textbf{Artifact Registry} \texttt{\{contrato\}}\desc{imagens da API, do hub e do agente}};
  \node[servicoger, alvo] (job) at (11.8,2.0)
    {\textbf{Cloud Run Job} \texttt{-migrate}\desc{antes do \emph{deploy} da API}};
  \node[servico, alvo] (run) at (11.8,0.0)
    {\textbf{Cloud Run} \texttt{-api}, \texttt{-frontend}\desc{nova revisão + sonda pós-\emph{deploy}}};
  \node[ia] (eng) at (11.8,-2.0)
    {\textbf{Agent Engine}\desc{só se \texttt{apps/agente-professor} mudou}};
  \node[balde] (exp) at (11.8,-4.0)
    {\textbf{Cloud Storage}\desc{\texttt{-exp} (\emph{bundles}), \texttt{-exp-mapas}}};
  \node[servico, alvo] (api) at (11.8,-6.0)
    {\textbf{API}\desc{\texttt{PUT} do manifesto: a \texttt{versaoAtiva} da experiência}};

  % barramento do token da DEPLOY_SA
  \draw[netpub] (gh.east) -- (wif.west) node[portlbl, midway, above=1pt]{OIDC};
  \draw[line width=0.9pt, draw=unooOrange] (wif.east) -- (7.4,0);
  \draw[line width=0.9pt, draw=unooOrange] (7.4,4.0) -- (7.4,-4.0);
  \foreach \alvo/\y/\rot in {ar/4.0/{\texttt{docker push}}, job/2.0/{\texttt{jobs execute}},
                              run/0.0/{\texttt{run deploy}}, eng/-2.0/{\texttt{PATCH}},
                              exp/-4.0/{\texttt{storage rsync}}}
    \draw[netpub] (7.4,\y) -- (\alvo.west) node[portlbl, midway, above=1pt]{\rot};

  % a ativação do manifesto não usa a DEPLOY_SA
  \draw[netpub] (gh.south) |- (api.west)
    node[portlbl, pos=0.75, above=1pt]{\emph{token} RS256 de 5 min};

  % de onde sai a imagem
  \draw[line width=0.8pt, draw=unooInkSoft, dashed, dash pattern=on 3pt off 2pt]
    (ar.east) -- (14.8,4.0) -- (14.8,-2.0);
  \coordinate (barra) at (14.8,0);
  \foreach \alvo in {job, run, eng}
    \draw[puxa] (barra |- \alvo.east) -- (\alvo.east);
  \node[portlbl, rotate=-90] at (14.8,1.0) {imagem};

\end{tikzpicture}}
\caption[Implantação detalhada --- o que o CI publica]{O que o CI publica, e
  onde. Um \emph{token} só, o da \texttt{DEPLOY\_SA}, trocado pelo WIF a cada
  execução (sem chave de longa duração), abre cinco portas: o \emph{push} das
  imagens, a migração, a revisão nova do Cloud Run, o \texttt{PATCH} da imagem
  do Agent Engine (pelo papel \texttt{agenteProfessorPublicador}) e o
  \texttt{rsync} dos \emph{bundles} e dos \emph{source maps}. A ativação da
  versão no manifesto é o único passo que fala com a API, e por isso usa a
  credencial que a API aceita. Os três serviços leem a imagem do mesmo Artifact
  Registry. O Agent Engine só nasce no segundo \emph{apply} do Terraform: antes
  dele, o \emph{job} do agente publica a imagem e para.}
\label{fig:deploy-publicacao}
\end{figure}

\section{Catálogo de recursos por contrato}

Este é o contrato formal com o repositório de Terraform. Nome divergente
significa CI quebrado.

\begin{tabularx}{\linewidth}{@{}L{30mm}L{44mm}X@{}}
\toprule
\headrow \thd{Recurso} & \thd{Nome esperado} & \thd{Consumido por}\\
\midrule
\endhead
Cloud Run (serviço) & \cd{\{contrato\}-frontend} (hub),
  \cd{\{contrato\}-api} & \arq{deploy.yml}\\
\zebra Cloud Run (\emph{job}) & \cd{\{contrato\}-migrate} --- na VPC, roda
  \cd{prisma migrate deploy} & \arq{deploy.yml}\\
Artifact Registry & \cd{\{contrato\}} (mesma região do Cloud Run) &
  \arq{deploy.yml}\\
\zebra \emph{Bucket} de experiências & \cd{\{contrato\}-exp} + Cloud CDN &
  \arq{deploy.yml}; hub (URL do \cd{iframe})\\
\emph{Bucket} de \emph{assets} & \cd{\{contrato\}-assets} + Cloud CDN, leitura
  pública & API (\env{ASSETS\_BUCKET}); mídia pesada das experiências\\
\zebra \emph{Bucket} de planos & \cd{\{contrato\}-planos} --- \textbf{privado},
  sem CDN, versionamento de objetos ligado & API (\env{PLANOS\_BUCKET})\\
Cloud SQL & \cd{\{contrato\}-pg} --- HA, PITR, \emph{pooler} ativo &
  API (\env{DATABASE\_URL})\\
\zebra Pub/Sub & \emph{topic} \cd{\{contrato\}-telemetria} + \emph{push sub}
  \cd{\{contrato\}-telemetria-ingest} \seta{} \rota{POST \{api\}/telemetry/ingest} &
  API\\
Secret Manager & \cd{\{contrato\}-database-url},
  \cd{\{contrato\}-db-password}, \cd{\{contrato\}-token-signing-key},
  \cd{\{contrato\}-auth-secret}, \cd{\{contrato\}-perfil-import-token},
  \cd{\{contrato\}-planos-import-token}, \cd{\{contrato\}-swagger-senha} e os
  três da Ajuda com avatar: \cd{\{contrato\}-liveavatar-api-key},
  \cd{\{contrato\}-ajuda-llm-segredo}, \cd{\{contrato\}-ajuda-sessao-chave} &
  \emph{env} do Cloud Run\\
\zebra Service Accounts & \cd{\{contrato\}-api} (Vertex, \emph{publisher} do
  Pub/Sub, \cd{objectAdmin} nos \emph{buckets} de \emph{assets} e de planos),
  \cd{\{contrato\}-run} (hub), \cd{\{contrato\}-push} (assina o \emph{push} da
  telemetria) & Cloud Run, Pub/Sub\\
Load Balancer & Um só, global: hub em \cd{\{hub\_domain\}}, API em
  \cd{\{api\_domain\}}, \emph{bundles} em \cd{\{exp\_domain\}}, \emph{assets} em
  \cd{\{assets\_domain\}} --- mesmo IP, \emph{host rules} & DNS\\
\zebra WIF para GitHub Actions & \emph{pool} \cd{github} &
  \arq{deploy.yml} (sem chave JSON)\\
\emph{Bucket} do agente & \cd{<projeto>-agente-professor} --- privado,
  versionado; prefixos \cd{bncc/}, \cd{livros/}, \cd{habilidades/},
  \cd{metadados/}, \cd{import-erros/} & \arq{apps/agente-professor/ingestao}
  (fora do Terraform); os \emph{data stores} leem dele\\
\zebra \emph{Data stores} (Vertex AI Search, \cd{global}) &
  \cd{agente-professor-habilidades} (estruturado, esquema próprio) e
  \cd{agente-professor-acervo} (\emph{parser} de \emph{layout},
  \emph{chunking} de 500 \emph{tokens} com títulos ancestrais) &
  o agente (\cd{discoveryengine.viewer})\\
Agent Engine & \cd{agente-professor} em \cd{us-east1} (dev e hml; produção a
  decidir). \textbf{Nasce só no segundo \emph{apply}}, quando
  \cd{agente\_professor\_imagem} aponta para uma imagem publicada &
  API (\env{AGENTE\_PROFESSOR\_ENGINE}); \arq{deploy.yml} (\cd{PATCH} da imagem)\\
\zebra SA do agente & \cd{\{prefixo\}-agente} (o id de SA tem teto de 30
  caracteres): \cd{aiplatform.user}, \cd{discoveryengine.viewer}, log, rastro,
  métrica e leitura do \emph{bucket} & o Agent Engine\\
Papéis custom no engine & \cd{agenteProfessorInvocador} (SA da API:
  \cd{reasoningEngines.get/query}, \cd{sessions.get/update}) e
  \cd{agenteProfessorPublicador} (CI: \cd{reasoningEngines.get/update}) &
  API; \arq{deploy.yml}\\
\bottomrule
\end{tabularx}

\begin{nota}
\textbf{Os recursos do agente de IA do professor} (\adrr{ADR-126},
\adrr{ADR-127}) vivem em \arq{agente\_professor.tf} no repositório de
fundação, que subiu o \emph{provider} \cd{google} para \cd{\textasciitilde> 7.46}
(o \cd{google\_vertex\_ai\_reasoning\_engine} com \cd{container\_spec} e o IAM do
engine não existem na 6.x). Três coisas ficam FORA do Terraform: a importação
dos documentos (\cd{documents:import}, pelo \arq{ingestao/} do \emph{monorepo}
--- dado não é infraestrutura), os \emph{service agents} do Discovery Engine e
do Agent Engine (criados uma vez por projeto com
\cd{gcloud beta services identity create}; o recurso que faria isso só existe
no \emph{provider} beta), e o \textbf{aumento de quota do Agent Engine}
(90 \cd{streamQuery} e 300 eventos de sessão por minuto, por projeto e região),
que se pede no console antes do \emph{go-live}. O \emph{chunking} do acervo nasce
certo ou não nasce: a configuração de processamento é imutável, e trocá-la é
recriar o \emph{data store} e reimportar tudo. Apagar o engine apaga o histórico
de todos os professores; o \cd{deletion\_policy} padrão recusa engine com
sessões, e não se troca para \cd{FORCE}.
\end{nota}

\begin{nota}
\textbf{A Ajuda com avatar não tem recurso de nuvem próprio} (\adrr{ADR-132}).
O \emph{vendor}, o LiveAvatar da HeyGen, fica fora do Google, e o que a
\emph{stack} ganha são três segredos no Secret Manager, todos com 32 ou mais
caracteres e lidos só pela API: \cd{\{contrato\}-liveavatar-api-key}
(\env{LIVEAVATAR\_API\_KEY}, a chave da conta, que nunca vai ao hub),
\cd{\{contrato\}-ajuda-llm-segredo} (\env{AJUDA\_LLM\_SEGREDO}, o \emph{Bearer}
que o \emph{vendor} manda ao nosso LLM; o MESMO valor fica guardado na conta do
\emph{vendor}) e \cd{\{contrato\}-ajuda-sessao-chave} (\env{AJUDA\_SESSAO\_CHAVE},
que assina o \emph{token} curto da sessão; só a API a conhece). O que fica
\textbf{fora do Terraform} são os ids do lado do \emph{vendor}: o contexto e a
configuração de LLM, que \arq{scripts/\B liveavatar-\B configurar.mjs} cria na
conta dele (idempotente pelo nome \cd{unooria-ajuda-<ambiente>}) e imprime como
\env{LIVEAVATAR\_CONTEXT\_ID} e \env{LIVEAVATAR\_LLM\_CONFIG\_ID} --- um
conjunto por contrato e ambiente, criado \emph{depois} de a API ter URL pública,
porque a configuração aponta para ela. Por isso a API sobe sem eles: faltou
qualquer peça, \rota{POST /ajuda/sessoes} responde 503 e o resto da plataforma
não percebe, e dá para aplicar o Terraform antes do \emph{script}. Trocar a URL ou
o segredo é apagar a configuração (e depois o segredo) no painel do
\emph{vendor} e rodar o \emph{script} de novo: o \emph{vendor} não tem edição
documentada, e apagar o segredo antes da configuração a quebra sem aviso.
\end{nota}

\begin{atencao}
\textbf{O mascote da Ajuda também não tem recurso próprio, mas pede duas coisas
ao caminho} (\adrr{ADR-158}). Ele usa o Vertex AI pelo mesmo ADC e a mesma SA da
API (o Gemini Live, na região \env{AJUDA\_MASCOTE\_REGIAO}, padrão
\cd{us-central1}: o Live não está no \cd{global}) e a mesma
\cd{\{contrato\}-ajuda-sessao-chave}; o modelo 3D
(\cd{ajuda/mascote/mascote-v1.glb}, 3,4\,MB) é um objeto do bucket de
\emph{assets}, publicado uma vez por contrato, atrás da CDN, nunca no
repositório. O que o caminho precisa:
\begin{enumerate}
  \item \textbf{\emph{Timeout} que cubra a sessão inteira.} A conversa é um
    WebSocket aberto do navegador à API (\rota{/ajuda/mascote/ao-vivo}) pelo
    mesmo LB e o mesmo Cloud Run. No Cloud Run, a conexão WebSocket morre no
    \emph{request timeout} do serviço; no LB com NEG \emph{serverless}, o
    \emph{timeout} de \emph{backend} é fixo em 60 minutos. Os 900\,s da API
    cobrem o padrão (10 minutos por sessão), mas o \emph{timeout} da API tem de
    ficar \textbf{acima de 10 minutos sempre}, e acompanhar o
    \cd{minutos-por-sessao} se o \cd{admin} subi-lo (o teto é 60) --- ou a
    conversa cai no meio com um fechamento que não é nosso. Sem afinidade de
    sessão: cada conversa é UMA conexão, e a instância que a recebe a leva até
    o fim.
  \item \textbf{O \emph{host} da API em \env{CORS\_ORIGINS} do lado certo}: a
    API confere o \cd{Origin} do \emph{upgrade} contra a lista do CORS; o hub
    precisa estar nela, como já está para o \cd{fetch}.
\end{enumerate}
\end{atencao}

\begin{decisao}
\textbf{Uma \emph{role} de banco, não duas.} O desenho original previa duas
\emph{roles}: uma dona do \emph{schema} para as migrações e uma \cd{app\_runtime}
sem \emph{ownership} para a aplicação. Na prática ficou \textbf{uma}, porque as
tabelas usam \cd{FORCE ROW LEVEL SECURITY} --- que sujeita o próprio dono às
políticas e torna a separação desnecessária. O que \emph{continua} obrigatório é
a \emph{role} não ter \cd{BYPASSRLS} nem \cd{SUPERUSER}; há SQL de verificação no
\emph{runbook} e um teste local (\arq{tests/tenancy.spec.ts}) que cobre o mesmo
contrato com uma \emph{role} sem \emph{ownership}.
\end{decisao}

\section{Caminhos de tráfego}

\begin{tabularx}{\linewidth}{@{}L{34mm}L{26mm}X@{}}
\toprule
\headrow \thd{Caminho} & \thd{Rede} & \thd{Observação}\\
\midrule
Navegador \seta{} LB \seta{} hub & Pública, TLS gerenciado &
  Sessão do hub; SSR e BFF.\\
\zebra Navegador \seta{} LB \seta{} CDN \seta{} \emph{bucket} de \emph{bundles} &
  Pública, TLS & O grosso dos \emph{bytes} do pico. Caminho versionado, portanto
  imutável; só o \arq{index.html} tem \cd{max-age=60}.\\
Navegador (dentro do \cd{iframe}) \seta{} LB \seta{} API & Pública, TLS +
  \emph{Bearer} & CORS por lista explícita (\env{CORS\_ORIGINS}).\\
\zebra hub \seta{} API & Pública, TLS &
  \env{API\_INTERNAL\_URL}; o hub repassa o \emph{token} do usuário.\\
API \seta{} Cloud SQL & \textbf{Privada}, VPC & O banco não tem IP público.\\
\zebra Cloud Run Job \seta{} Cloud SQL & \textbf{Privada}, VPC &
  Único caminho para rodar migração.\\
API \seta{} Vertex AI & Pública, TLS + ADC & Endpoint \cd{global}.\\
\zebra Navegador \seta{} LB \seta{} hub \seta{} API \seta{} Agent Engine &
  Pública, TLS; SSE em cada salto & O chat do agente (\adrr{ADR-128}). A API
  chama \cd{\{região\}-aiplatform.googleapis.com} com o ADC dela; teto de
  240\,s na API, 250\,s na rota de \emph{stream} do hub, 300\,s no Cloud Run do
  hub, 900\,s no da API e 60 minutos fixos
  no LB com NEG \emph{serverless}.\\
Agent Engine \seta{} Gemini e \emph{data stores} & Pública, TLS + SA do
  agente & Modelo em \cd{global} com reserva em \cd{us}; \emph{data stores} em
  \cd{global}. O navegador nunca chega aqui.\\
\zebra Navegador \seta{} LB \seta{} API (\rota{WS /ajuda/mascote/ao-vivo}) &
  Pública, WSS + bilhete & A conversa de voz do mascote da Ajuda
  (\adrr{ADR-158}), direto na URL pública da API, sem passar pelo hub (o Next
  não repassa WebSocket). \cd{Origin} conferido contra \env{CORS\_ORIGINS};
  bilhete de 2 minutos na primeira mensagem. Dura até o teto da sessão: o
  \emph{timeout} do Cloud Run da API tem de ser maior que ele.\\
API \seta{} Vertex AI (Gemini Live) & Pública, WSS + ADC & O
  \cd{BidiGenerateContent} em \cd{us-central1-aiplatform.googleapis.com}
  (\env{AJUDA\_MASCOTE\_REGIAO}). Só a API fala com ele; a instrução do sistema
  nunca vai ao navegador.\\
API \seta{} LiveAvatar (\cd{api.liveavatar.com}) & Pública, TLS +
  \cd{x-api-key} & Abrir e fechar a sala da Ajuda (\rota{/v1/sessions}), com
  \emph{timeout} de 10\,s (\adrr{ADR-132}). A chave nunca sai da API.\\
LiveAvatar \seta{} LB \seta{} API (\rota{/ajuda/llm}) & Pública, TLS +
  \emph{Bearer} + \emph{token} da sessão & O \emph{vendor} (na AWS;
  \cd{us-east-2} no \emph{spike}) chama \rota{POST /ajuda/llm/chat/completions} e
  \rota{GET /ajuda/llm/models} pela URL pública da API, o mesmo LB e o mesmo Cloud
  Run --- sem rota nova de rede. Não pode ficar atrás de IAP nem de \emph{allowlist}
  por IP; o Cloud Armor não pode recusar o corpo (JSON OpenAI com o \emph{prompt}
  do contexto); e a resposta é SSE com o primeiro \emph{byte} em menos de 10\,s,
  então nada de \emph{buffering} no caminho. Em desenvolvimento local, só com túnel
  (\cd{cloudflared}).\\
\zebra Navegador \seta{} LiveAvatar (LiveKit) & Pública, WebRTC & O vídeo da Ajuda
  não passa pela \emph{stack}. O \emph{firewall} da escola precisa liberar
  \cd{*.livekit.cloud} TCP 443 (WSS), \cd{*.turn.livekit.cloud} 443 (TURN/TLS),
  \cd{*.host.livekit.cloud} UDP 3478 e \cd{api.liveavatar.com} 443; recomendado
  UDP 50000--60000. Rede que bloqueia WebRTC mostra \enquote{não conseguimos
  conectar o vídeo} no modal.\\
API \seta{} Pub/Sub \emph{topic} & Pública, TLS + SA &
  Publicação do lote analítico.\\
Pub/Sub \seta{} API (\rota{/telemetry/ingest}) & Pública, TLS +
  \emph{ID token} & Verificado por \cd{PubSubPushGuard}: emissor Google,
  \cd{aud} e \cd{email} da SA de \emph{push}.\\
\zebra hub \seta{} WSO2 IS & Pública, TLS & OIDC. Fora da \emph{stack} do
  contrato.\\
\bottomrule
\end{tabularx}

\begin{atencao}
O \emph{Load Balancer} exige \textbf{dois} certificados gerenciados quando o hub
e a API estão em domínios distintos --- um por domínio. Provisionar apenas o do
hub faz a API responder com erro de TLS que \emph{parece} problema de aplicação:
o navegador não completa o \emph{handshake} e a experiência mostra falha de rede
sem nenhuma linha no log da API.
\end{atencao}

\section{Identidade de carga de trabalho e segredos}

\begin{figure}[H]
\centering
\fitwidth{%
\begin{tikzpicture}[x=1cm, y=1cm,
  passo/.style={etapa, text width=30mm, minimum height=15mm},
  origem/.style={etapaio, text width=30mm, minimum height=15mm},
  destino/.style={etapaout, text width=30mm, minimum height=15mm},
  fileira/.style={font=\scriptsize\bfseries\color{unooNavy}, anchor=south west, inner sep=1pt}]

  % quatro fileiras, uma por identidade; colunas em x = 0, 4.0, 8.3, 12.6
  \foreach \y/\rot in {0/{CI --- publicar}, -2.9/{API --- chamar a IA},
                       -5.8/{Agente --- consultar modelo e acervo},
                       -8.7/{Pub/Sub --- entregar a telemetria}}
    \node[fileira] at (-1.6,{\y+0.95}) {\rot};

  \node[origem] (gh) at (0,0)
    {\textbf{GitHub Actions}\desc{\emph{job} com \texttt{id-token: write}}};
  \node[passo] (oidc) at (4.0,0)
    {\emph{token} OIDC do GitHub\desc{repo, \emph{ref}, \emph{workflow}}};
  \node[passo, text width=35mm] (wif) at (8.3,0)
    {\textbf{WIF} \emph{pool} \texttt{github}\desc{valida e troca}};
  \node[destino] (dsa) at (12.6,0)
    {\emph{access token} da \texttt{DEPLOY\_SA}\desc{curta duração}};

  \node[origem] (run) at (0,-2.9)
    {\textbf{Cloud Run} \texttt{-api}\desc{sem \emph{env} de credencial}};
  \node[passo] (md) at (4.0,-2.9)
    {\emph{metadata server}\desc{ADC, nível 3}};
  \node[passo, text width=35mm] (saapi) at (8.3,-2.9)
    {SA \texttt{\{contrato\}-api}\desc{\texttt{aiplatform.user};\\\texttt{agenteProfessorInvocador}}};
  \node[destino] (vx) at (12.6,-2.9)
    {\textbf{Gemini} e \textbf{Agent Engine}\desc{autorizados}};

  \node[origem] (eng) at (0,-5.8)
    {\textbf{Agent Engine}\desc{o agente em execução}};
  \node[passo] (mdeng) at (4.0,-5.8)
    {credencial do \emph{runtime}\desc{a SA do engine}};
  \node[passo, text width=35mm] (saag) at (8.3,-5.8)
    {SA \texttt{\{prefixo\}-agente}\desc{\texttt{aiplatform.user};\\\texttt{discoveryengine.viewer}}};
  \node[destino] (ds) at (12.6,-5.8)
    {\textbf{Gemini} e os dois \emph{data stores}\desc{autorizados}};

  \node[origem] (sub) at (0,-8.7)
    {\emph{push sub}\\\texttt{-telemetria-ingest}};
  \node[passo] (idt) at (4.0,-8.7)
    {\emph{ID token} da SA \texttt{\{contrato\}-push}\desc{\texttt{oidc\_token} da \emph{subscription}}};
  \node[passo, text width=35mm] (guard) at (8.3,-8.7)
    {\texttt{PubSubPushGuard}\desc{emissor Google,\\\texttt{aud} e \texttt{email}}};
  \node[destino] (ing) at (12.6,-8.7)
    {\texttt{POST /telemetry/\B ingest}\desc{aceito}};

  \foreach \a/\b in {gh/oidc, oidc/wif, wif/dsa, run/md, md/saapi, saapi/vx,
                     eng/mdeng, mdeng/saag, saag/ds, sub/idt, idt/guard, guard/ing}
    \draw[pipe] (\a) -- (\b);

\end{tikzpicture}}
\caption[Identidade de carga de trabalho]{Os quatro caminhos de identidade de
  máquina, todos sem chave de longa duração: WIF para o CI, ADC para a API, a
  conta de serviço do próprio engine para o agente, e o \emph{ID token} assinado
  pelo Google para a entrega da telemetria. Nenhum arquivo JSON de conta de
  serviço existe em nenhum lugar do fluxo.}
\label{fig:identidade-maquina}
\end{figure}

\begin{nota}
\textbf{A mesma cadeia nos três ambientes.} A API resolve a credencial com
\cd{google-auth-library}, que tenta, na ordem: (1)
\env{GOOGLE\_APPLICATION\_CREDENTIALS}; (2) o arquivo de
\cd{gcloud auth application-default login}; (3) o \emph{metadata server}. Antes
disso havia \emph{só} o nível 3 --- e por isso nenhuma chamada de IA funcionava
fora do Cloud Run: nem no \emph{host}, nem no contêiner, onde
\cd{metadata.google.internal} não resolve. Em desenvolvimento a API ficava presa
em \env{AI\_PROVIDER=mock} e a app-90 devolvia um plano falso.
\end{nota}

\begin{atencao}
\textbf{Nunca chave de API.} Não há campo de \emph{env} para chave de IA em lugar
nenhum do projeto --- nem em \arq{.env.example}. Os dois \emph{compose} de
desenvolvimento montam o ADC do \emph{host} em \arq{/adc} somente leitura; em
Cloud Run a SA do serviço \emph{é} a credencial, e o que falta é apenas o papel
\cd{roles/aiplatform.user}. Sem ele, toda geração falha com 403 do Vertex --- e o
403 vem do Vertex, não da API, o que é um sintoma confusamente parecido com
quota esgotada.
\end{atencao}

\begin{tabularx}{\linewidth}{@{}L{46mm}L{20mm}X@{}}
\toprule
\headrow \thd{Segredo / variável} & \thd{Onde vive} & \thd{Por quê}\\
\midrule
\cd{\{contrato\}-database-url}, \cd{\{contrato\}-db-password} &
  Secret Manager & Credencial de banco.\\
\zebra \cd{\{contrato\}-token-signing-key} & Secret Manager &
  Chave privada RS256 que assina o \emph{token} de experiência. Também usada
  pelo CI para assinar o \emph{token} de cinco minutos do \emph{upsert} de
  manifesto.\\
\cd{\{contrato\}-auth-secret} & Secret Manager & Segredo de sessão do
  NextAuth.\\
\zebra \cd{\{contrato\}-perfil-import-token},
  \cd{\{contrato\}-planos-import-token} & Secret Manager &
  Chaves de serviço de duas rotas de importação. \textbf{Ausentes = rota
  desligada} (503), nunca aberta.\\
\env{OIDC\_ISSUER}, \env{OIDC\_JWKS\_URI} & \emph{env} em texto plano &
  \textbf{Não são segredo}: são URLs públicas do hub. Ficam em texto plano
  derivadas de \cd{hub\_domain} --- é isso que garante que o \cd{iss} do
  \emph{token} e o esperado pela API \emph{não podem} divergir.\\
\zebra \env{TELEMETRIA\_PUSH\_SA}, \env{TELEMETRIA\_PUSH\_AUDIENCE} &
  \emph{env} em texto plano & Identificam quem pode chamar
  \rota{/telemetry/ingest}. Ausentes = rota 401.\\
\env{AGENTE\_\B PROFESSOR\_\B MODO}, \env{AGENTE\_\B PROFESSOR\_\B ENGINE} &
  \emph{env} em texto plano & \cd{desligado} enquanto o engine não existe;
  \cd{agent-engine} e o \emph{resource name} completo depois do segundo
  \emph{apply}. Não são segredo: quem autoriza é o IAM do engine.\\
\zebra \cd{\{contrato\}-liveavatar-api-key}, \cd{\{contrato\}-ajuda-llm-segredo},
  \cd{\{contrato\}-ajuda-sessao-chave} & Secret Manager & A Ajuda com avatar
  (\adrr{ADR-132}): a chave da conta do \emph{vendor}, o \emph{Bearer} que ele
  manda ao nosso LLM e a chave que assina o \emph{token} da sessão. Ausentes =
  Ajuda desligada (503), nunca aberta; a API sobe normal.\\
\env{LIVEAVATAR\_\B CONTEXT\_\B ID}, \env{LIVEAVATAR\_\B LLM\_\B CONFIG\_\B ID},
  \env{LIVEAVATAR\_\B SANDBOX} & \emph{env} em texto plano & Os ids que
  \arq{scripts/\B liveavatar-\B configurar.mjs} imprime, fora do Terraform. Não
  são segredo: sem a chave da conta não abrem nada. \env{LIVEAVATAR\_SANDBOX=1}
  (o avatar de testes, \textasciitilde{}1 minuto, sem crédito) é recusado pelo Zod
  em produção.\\
\bottomrule
\end{tabularx}

\section{Dimensionamento}

\begin{tabularx}{\linewidth}{@{}L{30mm}L{34mm}X@{}}
\toprule
\headrow \thd{Recurso} & \thd{Configuração} & \thd{Motivo}\\
\midrule
Cloud Run --- API & CPU sempre alocada, \cd{min-instances} $\geq$ 1,
  concorrência 80 &
  Há trabalho \emph{in-process} fora do ciclo de requisição: o \emph{worker} de
  geração de conteúdo da exp08ef1b e a geração de plano de aula da app-90. Com CPU
  liberada entre requisições, a geração morreria no meio.\\
\zebra Cloud Run --- hub & \cd{min-instances} $\geq$ 1 &
  Elimina \emph{cold start} no pico de \emph{login}.\\
Cloud SQL & Instância pequena, HA, PITR, \emph{pooler} &
  \num{10000} estudantes e o volume medido de telemetria cabem folgadamente. O
  \emph{pooler} é obrigatório porque o \emph{autoscaling} do Cloud Run
  multiplica \emph{pools} do Prisma e estoura \cd{max\_connections}.\\
\zebra CDN & \emph{cache} longo no caminho versionado; \cd{max-age=60} no
  \arq{index.html} & O caminho carrega a versão, logo o conteúdo é imutável. O
  \arq{index.html} é o ponteiro: \emph{cache} longo nele atrasaria um
  \emph{rollback}.\\
Pub/Sub & \emph{Push subscription} com \cd{oidc\_token} &
  Amortecedor durável do pico, e uma peça a menos que um \emph{worker} dedicado.\\
\bottomrule
\end{tabularx}

\subsection{Volumes medidos}

\begin{tabularx}{\linewidth}{@{}L{40mm}L{28mm}X@{}}
\toprule
\headrow \thd{Grandeza} & \thd{Valor medido} & \thd{Como se chegou a ele}\\
\midrule
\emph{Payload} de uma partida da app-70 & \SI{1.1}{\kilo\byte} de JSON
  (\SI{1.5}{\kilo\byte} em JSONB) &
  Medido em 2026-08-26. O jogo tem sete pontos de log, todos atos discretos com
  guarda própria: 8 Chamas, 4 aldeões, 1 Borda + 3 Destaques, 4 cartões, 1
  primeira conversa --- teto real de \textasciitilde22 eventos por partida. Nada
  por quadro, nada por metro andado.\\
\zebra Armazenamento da app-70 por contrato & \textbf{\textasciitilde\SI{30}{\mega\byte}} &
  \num{10000} estudantes jogando uma vez. A estimativa \emph{anterior} --- centenas
  de MB a \SI{2.4}{\giga\byte} --- partia do teto do contrato (2000 eventos) e
  errava por duas ordens de grandeza. O teto do contrato caiu de 2000 para 200,
  que é o que de fato limita a linha se um cliente adulterado mandar lixo.\\
Custo de gravação da telemetria & \textasciitilde16 \cd{INSERT}/s no pico
  (era \textasciitilde800) &
  Um lote de 50 eventos custava 50 \emph{round-trips}; em \cd{INSERT}
  multi-linha custa um. No pico de \textasciitilde4000 simultâneos a diferença é
  entre saturar e não notar.\\
\zebra \emph{Prompt} de um plano de aula anual & \textasciitilde\num{190000}
  \emph{tokens} &
  Currículo da rede + as quatro unidades do ano. É o que justifica a quota
  separada de professor (\num{2000000}/dia $\approx$ 6 a 9 planos).\\
Mídia da exp04ef1b & \SI{52}{\mega\byte} (21 arquivos MP4) &
  Fora do \emph{git} e fora do \emph{bundle}: vai para o \emph{bucket} de
  \emph{assets}, referenciada por \env{VITE\_VIDEOS\_BASE\_URL}.\\
\bottomrule
\end{tabularx}

\section{Ambientes}

\begin{tabularx}{\linewidth}{@{}L{32mm}L{34mm}X@{}}
\toprule
\headrow \thd{Ambiente} & \thd{Como sobe} & \thd{Diferenças relevantes}\\
\midrule
\emph{Host} (desenvolvimento) & \cd{docker compose up -d pg} +
  \cd{pnpm dev} &
  Tudo no \emph{host}; só o banco em contêiner. \env{AI\_PROVIDER} pode ser
  \cd{mock}; \emph{drivers} de armazenamento locais; sem Pub/Sub.\\
\zebra Contêiner (desenvolvimento) & \cd{pnpm docker:dev} &
  A \emph{stack} inteira em contêiner --- banco, API, hub e as cinco
  experiências --- com o código do \emph{host} por \emph{bind mount}: salvar
  recompila lá dentro. Mesmas portas do \cd{pnpm dev}, então é um \emph{ou} o
  outro.\\
Imagens congeladas & \cd{docker compose up -d --build} &
  Aproxima-se da produção: imagens construídas, migração e \emph{seed}
  automáticos no primeiro \cd{up}.\\
\zebra \cd{dev} e \cd{hml} em GCP & CI no \emph{push} para \cd{dev} ou \cd{hml} &
  \emph{Stack} real do piloto (\cd{plataforma-unooria-dev}), com todos os
  recursos da Figura~\ref{fig:deploy-detalhe}.\\
Produção & Mesmo CI (\cd{main} entra no gatilho quando virar produção), outras \cd{vars} &
  O \emph{workflow} não conhece o nome de nenhum recurso: tudo entra por
  \cd{vars}, cujos valores saem de \cd{terraform output}.\\
\bottomrule
\end{tabularx}

\begin{nota}
Nos ambientes de contêiner, a leitura \emph{server-side} do
\arq{experiencia.json} exige \env{EXP\_INTERNAL\_HOST} --- que troca apenas o
\emph{host} e \textbf{preserva a porta}. \env{EXP\_INTERNAL\_URL}, que é um
endereço fixo, só serve quando existe \emph{uma} origem de \emph{bundles} (a CDN,
em produção): em desenvolvimento, cada experiência tem um \emph{dev server} numa
porta diferente, e um endereço fixo faria todas lerem o manifesto de uma só.
O sintoma dessa configuração errada é específico e memorável: o UnooTube, que não
declara \cd{viewport}, aparece dentro da casca de celular da Instanooria.
\end{nota}

\begin{atencao}
Os \emph{dev servers} das experiências declaram
\cd{allowedHosts: ["host.docker.internal"]}. O Vite 5.4+ recusa cabeçalho
\cd{Host} que não conheça, e sem essa declaração o hub em contêiner não lê o
manifesto --- e a experiência aparece esticada na janela em vez de dentro do
aparelho.
\end{atencao}

\section{O provedor de identidade fora da \emph{stack}}

O WSO2 Identity Server foge do padrão \cd{\{contrato\}-*} por desenho: é
\textbf{único para vários contratos} e tem ciclo de vida próprio. Vive em projeto
GCP próprio e, em desenvolvimento, num \emph{compose} separado
(\arq{wso2/docker-compose.yml}, IS em \cd{:9443} e banco próprio em
\cd{:5433}).

\begin{decisao}
\textbf{O hub assina os \emph{tokens} de experiência enquanto o IdP central não
entra} (\adrr{ADR-15}). O hub emite RS256 com TTL de 15 minutos e publica o JWKS
em \rota{/.well-known/jwks.json}; a API valida por \env{OIDC\_ISSUER} e
\env{OIDC\_JWKS\_URI}. Trocar para o WSO2 é trocar essas duas variáveis --- o
\cd{OidcGuard} não muda uma linha. A alternativa (esperar o IdP para começar a
integrar) faria toda a plataforma depender de um servidor Java em
\emph{bootstrap} para que qualquer experiência pudesse ser testada.
\end{decisao}

\begin{atencao}
O \cd{issuer} do WSO2 é \emph{tenant-qualified}
(\cd{.../t/\{contrato\}/oauth2/token}) e deve ser copiado \textbf{literalmente}
do \emph{discovery}, nunca montado por concatenação. Um \cd{iss} montado à mão
que difere por uma barra faz toda validação falhar com
\enquote{\emph{issuer} inválido} --- e a mensagem não diz qual dos dois lados
está errado.
\end{atencao}

\section{Observabilidade, continuidade e reversão}

\begin{tabularx}{\linewidth}{@{}L{34mm}X@{}}
\toprule
\headrow \thd{Aspecto} & \thd{Como é atendido hoje}\\
\midrule
Sanidade de processo & \rota{GET /healthz} (processo vivo) e
  \rota{GET /readyz} (dependências prontas). São rotas sem \emph{guard}, de
  propósito.\\
\zebra Log & Cloud Logging pelo \emph{stdout} do Cloud Run. \textbf{Nenhuma PII
  em texto claro}; texto livre de bem-estar nunca vai a log (só contagens).\\
Métricas & Métricas nativas do Cloud Run e do Cloud SQL. Não há métrica de
  aplicação customizada --- é lacuna conhecida (Cap.~\ref{cap:riscos}).\\
\zebra Reversão do hub e da API & Revisões do Cloud Run. \cd{gcloud run
  services update-traffic} devolve o tráfego à revisão anterior.\\
Reversão de experiência & \cd{node apps/hub/scripts/publicar-manifesto.mjs
  <exp> --versao <anterior>}. Sem \emph{build}, sem imagem.\\
\zebra Reversão de conteúdo de professor & Versionamento de objetos do
  \emph{bucket} de planos: cada \enquote{Salvar} é uma versão, sem uma linha de
  código de histórico.\\
Recuperação do banco & HA + PITR do Cloud SQL.\\
\zebra Migração de \emph{schema} & Cloud Run Job dentro da VPC, executado
  \textbf{antes} do \emph{deploy} do serviço, com \cd{--wait} para que a falha
  interrompa o \emph{pipeline}.\\
\bottomrule
\end{tabularx}

\begin{limitacao}
Não há \emph{alerting} configurado, nem painel. Há log e métrica nativa, e há
sanidade por \rota{/readyz} --- o que basta para diagnosticar, mas não para ser
avisado; e \rota{/readyz} não é chamada por ninguém, porque o \emph{deploy}
não configura sonda. O estudo completo --- o que se vê hoje, elo por elo, os
sinais a acrescentar, a comparação de ferramentas e o roteiro em fases com
custo por contrato --- é o Capítulo~\ref{cap:observabilidade}. A fase 0 dele
(canais, \emph{uptime checks}, políticas, SLOs e painel, tudo em Terraform) é
o que fecha o risco \cd{R-05} antes de novembro de 2026.
\end{limitacao}
